\begin{abstract}
\noindent
This document provides a comprehensive, progressive guide to multi-agent systems architectures, beginning with foundational concepts of large language models and advancing through increasingly complex architectures. Starting from single-agent systems, we progressively introduce slave hierarchies, orchestration patterns, and multi-layered hierarchies. The guide explores three agent execution models (streaming, batch, and parallel), multiple communication patterns (direct, file-mediated, and hybrid), and diverse control models (imperative, declarative, and hybrid). We examine orchestration cycles ranging from single-pass to iterative and reactive patterns, alongside system topologies including serial, parallel, and tree configurations. Central to this framework are four canonical patterns: pipeline, fan-out/fan-in, master-slave, and feedback loop architectures. Throughout, we provide practical guidance for selecting and combining patterns based on specific task characteristics, computational constraints, and system requirements.
\end{abstract}
